\section{Results}
\label{sec:results}

We use \benchname{} to ask three questions. Does \fmname{} provide
forecast signal that persistence does not (\texttt{b1},
\texttt{b3}, \texttt{b6})? Does routing the LLM through structured
evidence improve supervisory tickets once precision, recall,
explanation quality, and false interventions are scored jointly
(\texttt{b4}, \texttt{b5})? Which components are load-bearing, and
which are reserves under degradation (ablations)?

\subsection{What the forecaster changes}
Before the aggregate scores, Table~\ref{tab:qualitative} shows the
mechanism on a single week. Given the same network, the raw-snapshot
LLM (\texttt{m2}) reads only present-time concentration and is
structurally capped at \texttt{Investigate}; the \fmname{} stress
scenarios let the LLM (\texttt{m7}) escalate a specific hub with a
quantified cascade as warrant. The same grounding can also misfire:
\S\ref{app:qualitative} pairs this success with a week where an
equally well-grounded escalation targets the wrong hubs (judge-confirmed
grounding, zero target precision)---the failure mode the aggregate FIR
quantifies below.

\begin{table}[t]
\centering\small
\setlength{\tabcolsep}{4pt}
\begin{tabularx}{\linewidth}{@{}l>{\raggedright\arraybackslash}X@{}}
\toprule
\multicolumn{2}{@{}>{\raggedright\arraybackslash}p{\linewidth}@{}}{%
\textbf{Same hub (node~2786, $9.35$B exposure), week 2025-03-24}, with
vs.\ without the forecaster. Opus-4.8 judge: \texttt{m2}~$=2$,
\texttt{m7}~$=4$.}\\
\midrule
\texttt{m2} & (no FM) \emph{moderate} risk; node~2786
\textsc{Investigate} ($0.62$): ``\dots without forward-looking signals
\textbf{no protocol meets the $0.75$ threshold for intervention}---warranting
investigation of the largest hubs.'' \\
\addlinespace
\texttt{m7} & (FM) \emph{elevated} risk; node~2786
\textsc{Recommend-Reduce} ($0.88$): ``dominant in \textbf{S1
single-failure cascade \dots\ $18.67\%$ system loss, $84$ distressed} \dots\
and largest contributor to S5 correlated stress.'' \\
\bottomrule
\end{tabularx}
\caption{A single week, with and without forecaster grounding (verbatim
ticket excerpts; one of three consistency runs). Without a forward signal
the LLM is structurally capped at \textsc{Investigate}; the \fmname{}'s
stress scenarios license a quantified escalation. The same grounding can
equally warrant a confident-but-wrong escalation---the over-intervention
case (2025-04-28) and full tickets are in \S\ref{app:qualitative}.}
\label{tab:qualitative}
\end{table}

\subsection{Comparison}
\paragraph{Forecast quality.}
On static error metrics, persistence is a deceptively strong
baseline: at $h\!=\!4$ it beats the FM on PageRank MAE
($3.4\!\times\!10^{-5}$ vs $4.5\!\times\!10^{-5}$) and edges it on
Spearman rank correlation ($0.570$ vs $0.558$). The FM contributes signal where
persistence is structurally pinned: trend consistency
$0.525$--$0.632$ across horizons (persistence: $0$);
\texttt{b3\_calibration} PI coverage $0.913$ against a $0.90$
target and ECE $0.013$; \texttt{b6\_robustness} mean relative
degradation $0.113$ vs $0.148$ for persistence ($24\%$ less). The
full per-horizon and per-regime audits are in
\S\ref{app:full-tables}.

\paragraph{Main result: decision quality.}
Table~\ref{tab:headline} consolidates the central
\texttt{b5\_decision} comparison. The persistence$+$rules baseline
(\texttt{m1}) achieves the highest single-axis precision
($0.720$); like \texttt{m2}/\texttt{m5} it never escalates beyond
\texttt{Investigate}, so its FIR is undefined, not a safety
guarantee. Against this baseline, the full
\fmname$+$LLM$+$gate stack (\texttt{m7}) lifts ticket F1 from
$0.0076$ to $0.0233$ ($+207\%$; $p\!<\!10^{-4}$). Against the comparable pure-LLM
baseline that consumes raw snapshots without the forecaster
(\texttt{m2}, the only other LLM-bearing row), \texttt{m7} lifts
LLM-as-judge explanation quality from $2.24$ to $2.45$ on a 1--5
scale ($+9.4\%$) under a tier-above Claude~Opus~4.8 judge
\citep{Zheng2023LLMJudge}; the lift is significant under GPT-5.5
($p\!<\!0.001$) but not Opus~4.8 ($p\!=\!0.23$;
\S\ref{app:stats}), so we read it as directional rather than
confirmed. The
cost surfaces sharply once \texttt{m2} is in the table: the
raw-snapshot LLM issues \emph{no} interventions, but adding the FM
signal (\texttt{m2}$\!\to\!$\texttt{m6}) makes it escalate, and
$44.8\%$ of those interventions misfire
($\mathrm{FIR}~0\!\to\!0.448$; $p\!<\!10^{-4}$). The safety gate
(\texttt{m6}$\!\to\!$\texttt{m7}) reduces FIR only marginally
($0.448\!\to\!0.437$), and its effects on F1 and judge score are
within noise (\S\ref{app:stats}). Both LLM variants
reach grounding $1.000$, so the failure mode is not citation
hallucination but \emph{over-reading the forecaster's evidence as
warrant for high-severity action}.\footnote{\texttt{m3\_evolvegcn}
and \texttt{m4\_fm\_only} are predictor-only and do not emit
tickets (forecast quality in \S\ref{app:b1-full}); EvolveGCN
trails persistence downstream, consistent with its $\sim$220
weekly-snapshot training budget.}

\paragraph{Judge robustness.}
We re-judge \texttt{m2}/\texttt{m6}/\texttt{m7} with a
cross-family panel (Table~\ref{tab:judge-panel}). All three
frontier judges rank \texttt{m7} highest (mean score varying by
$\leq\!0.45$). The one disagreement, where Gemini~2.5~Pro scores
\texttt{m2} above \texttt{m6} ($2.69$ vs $2.55$, still below
\texttt{m7}), sits within Gemini's ${\sim}3\times$ higher weekly
noise (std $1.3$--$1.5$ vs $0.3$--$0.6$). Cross-family agreement
on the top system closes the judge-family confound.

\begin{table*}[t]
\centering
\small
\setlength{\tabcolsep}{4.5pt}
\begin{tabular*}{\linewidth}{@{\extracolsep{\fill}}lllccccccc@{}}
\toprule
ID & Pred. & Decision & Prec.$\uparrow$ & Recall$\uparrow$ & F1$\uparrow$ & Judge$\uparrow$ & FIR$\downarrow$ & Ground.$\uparrow$ & Stabil.$\uparrow$ \\
\midrule
\texttt{m1} & persist.   & rules               & \textbf{0.720} & 0.004 & 0.0076          & ---            & ---            & ---   & 0.514 \\
\texttt{m2} & ---        & LLM (no FM)         & 0.575          & 0.009 & 0.0183          & 2.24           & ---            & 1.000 & 0.532 \\
\texttt{m5} & \fmname{}  & rules               & 0.600          & 0.010 & 0.0190          & ---            & ---            & ---   & 0.257 \\
\texttt{m6} & \fmname{}  & LLM                 & 0.570          & 0.012 & \textbf{0.0241} & 2.41           & 0.448          & 1.000 & 0.488 \\
\texttt{m7} & \fmname{}  & LLM + safety gate   & 0.580          & 0.012 & 0.0233          & \textbf{2.45}  & 0.437          & 1.000 & 0.435 \\
\bottomrule
\end{tabular*}
\caption{Headline \texttt{b5\_decision} comparison on the frozen
2025 test split ($n\!=\!29$). Rows are read as a multi-objective
trade-off rather than a single leaderboard: \texttt{m1} is the
conservative high-precision baseline; \texttt{m2} is the pure-LLM
baseline (no forecaster); \texttt{m6}/\texttt{m7} are the
FM-grounded LLM variants. Judge scores (1--5 scale) are from a
Claude~Opus~4.8 judge, strictly tier-above the Claude~Opus~4.7
decision LLM; cross-family agreement is in
Table~\ref{tab:judge-panel}. A dash in the FIR column means
undefined: \texttt{m1}/\texttt{m2}/\texttt{m5} emit no
intervention-level ticket that could misfire.}
\label{tab:headline}
\end{table*}

\begin{table}[t]
\centering\small
\setlength{\tabcolsep}{4pt}
\begin{tabular}{@{}lcccc@{}}
\toprule
Method & Opus~4.8 & Gemini~2.5~Pro & GPT-5.5 & Haiku~4.5\textsuperscript{\dag} \\
\midrule
\texttt{m2} & 2.24 & 2.69 & 1.90 & 2.03 \\
\texttt{m6} & 2.41 & 2.55 & 2.45 & 2.48 \\
\texttt{m7} & \textbf{2.45} & \textbf{2.90} & 2.45 & 2.69 \\
\bottomrule
\end{tabular}
\caption{Cross-family judge panel on the same Opus~4.7 decisions
($n\!=\!29$). Gemini and GPT use $\mathrm{reasoning}\!=\!
\{\mathrm{effort:minimal}\}$ so the thinking budget matches the
Anthropic baseline. \textsuperscript{\dag}~Haiku~4.5 is a
tier-below weak-judge reference.}
\label{tab:judge-panel}
\end{table}

\paragraph{Decision-LLM swap.}
The headline operating point is not Opus-bound: a $\sim$5$\times$
cheaper decision LLM Pareto-dominates it. We re-run the decision
layer of \texttt{m6} and \texttt{m7} on the \emph{same} FM-derived
prompts with two non-Opus~4.7 decision LLMs (Claude~Sonnet~4.6,
within-family and tier-below; and Google~Gemini~2.5~Pro,
cross-family), and re-score with the same Opus~4.8 judge
(Table~\ref{tab:decision-panel}).
On both architectures Sonnet~4.6 is a Pareto improvement over
Opus~4.7: F1 rises ($+15\%$/$+24\%$; $p\!<\!0.001$ on
\texttt{m7}), judge quality is directionally higher
($+0.11$/$+0.21$; n.s.), and FIR collapses from $0.437$--$0.448$
to $0.000$, and genuinely so: Sonnet~4.6 still issues interventions at
Opus-comparable volume ($\sim$190 vs $252$), all on truly-stressed
targets. Gemini~2.5~Pro is more conservative still (lower F1,
FIR~$0$), so over-intervention is a property of the decision model,
not of LLM decision layers per se. The
operating point a practitioner would now deploy is \texttt{m7}
with Sonnet~4.6 (F1 $0.0288$, FIR $0.000$, judge $2.66$).

\begin{table}[t]
\centering\small
\setlength{\tabcolsep}{4pt}
\begin{tabular}{@{}llccc@{}}
\toprule
Method & Decision LLM & F1$\uparrow$ & FIR$\downarrow$ & Judge$\uparrow$ \\
\midrule
\texttt{m6} & Opus~4.7 (main)  & 0.0241 & 0.448 & 2.41 \\
\texttt{m6} & Sonnet~4.6       & \textbf{0.0276} & \textbf{0.000} & \textbf{2.52} \\
\texttt{m6} & Gemini~2.5~Pro   & 0.0129 & 0.000 & 2.21 \\
\midrule
\texttt{m7} & Opus~4.7 (main)  & 0.0233 & 0.437 & 2.45 \\
\texttt{m7} & Sonnet~4.6       & \textbf{0.0288} & \textbf{0.000} & \textbf{2.66} \\
\texttt{m7} & Gemini~2.5~Pro   & 0.0139 & 0.000 & 2.38 \\
\bottomrule
\end{tabular}
\caption{Cross-decision-LLM panel on the same FM-derived prompts,
judge fixed at Claude~Opus~4.8 ($n\!=\!29$). Sonnet~4.6 dominates
Opus~4.7 on both architectures. The bold row
(\texttt{m7}~$\times$~Sonnet~4.6) beats every row in
Table~\ref{tab:headline} on every LLM-bearing axis at
$\sim$5$\times$ lower cost.}
\label{tab:decision-panel}
\end{table}

\paragraph{Layer-wise contribution.}
Holding the decision policy fixed at \emph{LLM} and adding the
\fmname{} signal lifts ticket F1 by $32\%$
($0.0183\!\to\!0.0241$; $p\!<\!0.001$); holding the predictor
fixed at \fmname{} and swapping rules for the LLM lifts F1 by
another $27\%$ ($0.0190\!\to\!0.0241$). The two layers therefore
make additive, non-substitutable contributions
(Figure~\ref{fig:layerwise} in \S\ref{app:full-tables}). Adding
the safety gate trades $3\%$ F1 for judge changes within noise;
its value is operational: every intervention clears two gates
with auditable logs.

\paragraph{Matched-budget analysis.}
F1 conflates target quality with volume:
\texttt{m6}/\texttt{m7} emit $6.3$--$6.5$ targets per week vs
$4.9$ for \texttt{m2} and $2.1$ for \texttt{m1}. Truncating every
method to its top-$k$ targets per week ($k\!\in\!\{1,2,3\}$;
\S\ref{app:stats-budget}) separates the two effects: the LLM
variants still roughly double the rules baseline's precision@$k$
($p\!\leq\!0.0055$), but the FM-fed and raw-snapshot LLM variants
become statistically indistinguishable. The FM's contribution is
thus not a sharper top of the ranking but support for a
\emph{larger} target set at undiminished precision; under
Opus~4.7 the expanded tail is where over-intervention
concentrates.

\subsection{Ablation}
\paragraph{Component ablation.}
Under clean 2025 data the confidence gate is load-bearing: removing
it raises FIR from $0.000$ to $0.429$. The scenario engine is
load-bearing for target extraction: skipping it collapses ticket
precision from $0.600$ to $0.000$. The data-health gate and
multi-horizon diversity are dormant on clean data but activate as
reserves under degradation. Under severe feature/edge masking
(80--98\%) the data-health gate suppresses up to $60\%$ of
interventions that would otherwise produce false alarms (FIR range
$[0.27,0.60]$ when disabled vs $0.000$ when triggered);
multi-horizon monitoring generates $\sim$4$\times$ more alerts in
the three crisis windows at unchanged precision.
Per-regime ablation audits are in \S\ref{app:ablation-stress}; the
historical SVB/USDC/Terra event study (lead times, not a test-split
result) is in \S\ref{app:b2-budget}.
